\subsection{乘积关于通常运算的行为}

命题 5. --- 设 \((\lambda_{\alpha})_{\alpha\in\mathbb{A}}\) 是 T 上的一个正测度族，关于关系 \(\leqslant\) 有向，在 \(\mathcal{M}(\mathrm{T})\) 中有一个上确界 \(\lambda\) 。为使正数值函数 g 局部 \(\lambda\) -可积，必须且只须 g 是局部 \(\lambda_{\alpha}\) -可积的（对每个 \(\alpha\in A\) ）且族 \((g\cdot\lambda_{\alpha})\) 在 \(\mathcal{M}(\mathrm{T})\) 中有上界；此时

\[
g \cdot \lambda = \sup _ {\alpha \in \mathrm{A}} g \cdot \lambda_ {\alpha}.
\]

条件的必要性是明显的。反之，设 \(g\) 对每个测度 \(\lambda_{\alpha}\) 都局部可积且族 \((g \cdot \lambda_{\alpha})_{\alpha \in \mathbb{A}}\) 有上界；用 \(\lambda'\) 表示它的上确界。函数 \(g\) 于是是 \(\lambda\) -可测的（§1, No.~4, 命题 11 的推论 2）；此外，对每个函数 \(h \in \mathcal{K}_{+}(\mathbf{T})\) ，

\[
\int^ {\bullet} (h g) d \lambda = \sup _ {\alpha \in A} \int^ {\bullet} (h g) d \lambda_ {\alpha} = \sup _ {\alpha \in A} \int^ {\bullet} h d (g \cdot \lambda_ {\alpha}) = \int^ {\bullet} h d \lambda^ {\prime}
\]

（§1, No.~4, 命题 11）。这首先蕴涵第一个端对任何 \(h\) 都是有限的，从而 \(g\) 局部 \(\lambda\) -可积；因此符号 \(\int^{\bullet}\) 可以换成 \(\int\) ，公式可以写成 \(\int h d(g \cdot \lambda) = \int h d\lambda'\) 。由此得出 \(g \cdot \lambda = \lambda'\) ，证明就此完成。

推论. --- 假设 \(\mu\) 是 T 上一个测度族 \((\mu_{\alpha})_{\alpha\in\mathrm{A}}\) 的和。为使定义在 T 上的正数值函数 g 局部 \(\mu\) -可积，必须且只须 g 是局部 \(\mu_{\alpha}\) -可积的（对每个 \(\alpha\in A\) ）且族 \((g\cdot\mu_{\alpha})_{\alpha\in\mathrm{A}}\) 可和。此时

\[
g \cdot \mu = \sum_ {\alpha \in A} g \cdot \mu_ {\alpha}.\tag{7}
\]

设 \((g_{\alpha})_{\alpha \in A}\) 是定义在 T 上的一个 \(\mu\) -可测正函数族。设 \(S_{\alpha}\) 是由 \(t \in T\) 中使 \(g_{\alpha}(t) \neq 0\) 者组成的集。我们称族 \((g_{\alpha})\) 是局部可数的，如果族 \((S_{\alpha})\) 是局部可数的（Ch. IV, §5, No.~9）；这等价于说，对 T 中每个紧集 K，由 \(\alpha \in A\) （使 \(g_{\alpha}|K\) 不为零者）组成的集是可数的。

命题 6. --- 设 \((g_{\alpha})_{\alpha \in A}\) 是定义在 T 上的局部 \(\mu\) -可积正数值函数的一个局部可数族。为使函数 \(g = \sum_{\alpha \in A} g_{\alpha}\) 局部 \(\mu\) -可积，必须且只须测度族 \((g_{\alpha} \cdot \mu)_{\alpha \in A}\) 可和，此时

\[
g \cdot \mu = \sum_ {\alpha \in A} g _ {\alpha} \cdot \mu .\tag{8}
\]

明显地，\(g\) 是 \(\mu\) -可测的（Ch. IV, §5, No.~2, 命题 4 及 No.~4, 定理 2 的推论 1）。因此，为使 \(g\) 局部 \(\mu\) -可积，必须且只须 \(\mu^{\bullet}(gf)\) 对每个 \(f \in \mathcal{K}_{+}(\mathrm{T})\) 都是有限的。现在，由于由 \(\alpha \in \mathrm{A}\) （使 \(g_{\alpha}f \neq 0\) 者）组成的集是可数的，我们有 \(\mu^{\bullet}(gf) = \sum_{\alpha \in \mathrm{A}} \mu^{\bullet}(g_{\alpha}f)\) （§1, No.~1, 命题 2 的推论）。置 \(\nu_{\alpha} = g_{\alpha} \cdot \mu\) ；条件 \(\mu^{\bullet}(gf) < +\infty\) 等价于条件 \(\sum_{\alpha \in \mathrm{A}} \nu_{\alpha}(f) < +\infty\) ：换言之，\(g\) 局部 \(\mu\) -可积当且仅当族 \((\nu_{\alpha})\) 可和。用 \(\nu\) 表示这个族的和，前面的计算给出等式 \(\nu(f) = \mu^{\bullet}(gf)\) ，它等价于 (8)。

推论. --- 设 \((g_{\alpha})\) 是局部 \(\mu\) -可积数值函数的一个序列，使得测度序列 \(g_{n} \cdot \mu\) 递增。为使这个序列在 T 上实测度所成的序向量空间 \(\mathcal{M}(T)\) 中有上界，必须且只须函数 \(g = \sup g_{n}\) 局部 \(\mu\) -可积；此时在 \(\mathcal{M}(T)\) 中序列 \((g_{n} \cdot \mu)\) 的上确界是测度 \(g \cdot \mu\) 。

只须把命题 6 应用于函数（局部几乎处处为正的）\(g_n' = g_{n+1} - g_n\) 。

命题 7. --- 设 X 是一个在无穷远处可数的局部紧空间，并设 \(t \mapsto \lambda_t\) 是一个 \(\mu\) -适宜映射（T 到 \(\mathcal{M}_+(X)\) 中的）。

设 \(g\) 是定义在 \(X\) 上的一个正数值函数，对测度 \(\nu = \int \lambda_t d\mu(t)\) 局部可积。于是由 \(t \in T\) 中使 \(g\) 不是局部 \(\lambda_t\) -可积者组成的集对 \(\mu\) 局部可忽略，映射 \(t \mapsto g \cdot \lambda_t\) （局部 \(\mu\) -几乎处处有定义）是 \(\mu\) -适宜的，且

\[
g \cdot \nu = \int (g \cdot \lambda_ {t}) d \mu (t).\tag{9}
\]

设 \((\mathrm{K}_{n})_{n\in\mathbf{N}}\) 是 X 的紧子集的一个递增序列，它们的内部覆盖 X；若 \(\eta\) 是 X 上任一正测度，则说 g 局部 \(\eta\) -可积等价于说 \(g\varphi_{\mathrm{K}_{n}}\) 对每个 n 都是 \(\eta\) -可积的。现在设 \(\mathrm{H}_{n}\) 是由 \(t\in\mathrm{T}\) 中使 \(g\varphi_{\mathrm{K}_{n}}\) 不是 \(\lambda_{t}\) -可积者组成的集，并设 \(\mathrm{H}=\bigcup_{n}\mathrm{H}_{n}\) ；由于 \(\mathrm{H}_{n}\) 对所有 n 都局部 \(\mu\) -可忽略（§3, No.~3, 定理 1），H 亦然，这就证明了陈述的第一个断言。把 \(\lambda_{t}\) 对 H 中的 t 换成 0（这不改变测度 \(\nu\) ），我们就可以假设 g 是局部 \(\lambda_{t}\) -可积的（对每个 \(t\in\mathrm{T}\) ）。对每个定义在 X 上的 \(\nu\) -可测正函数 h，我们依命题 3 及 §3, No.~2 的命题 5 有

\[
\int^ {\bullet} h d (g \cdot \nu) = \int^ {\bullet} (g h) d \nu = \int^ {\bullet} d \mu (t) \int^ {\bullet} (g h) d \lambda_ {t} = \int^ {\bullet} d \mu (t) \int^ {\bullet} h d (g \cdot \lambda_ {t}).
\]

这个公式与 §3, No.~2 的命题 5 首先表明（取 \(h \in \mathcal{K}_{+}(\mathbf{X})\) ）：映射 \(t \mapsto g \cdot \lambda_{t}\) 是标量本质 \(\mu\) -可积的，且其积分为 \(g \cdot \nu\) ；换言之，关系式 (9) 成立。其次，把 \(\mu\) 换成一个正测度 \(\mu' \leqslant \mu\) ，并取 \(h\) 为一个正的下半连续函数：由这些关系式立即得出 \(t \mapsto g \cdot \lambda_{t}\) 是 \(\mu\) -适宜的（§3, No.~1, 定义 1）。

命题 8. --- 设 \(\theta\) 是 T 上一个复测度，\(g_{1}\) 是一个局部 \(\theta\) -可积的复函数，\(\theta_{1}\) 是测度 \(g_{1} \cdot \theta\) 。为使定义在 T 上的复函数 \(g_{2}\) 局部 \(\theta_{1}\) -可积，必须且只须 \(g_{2} g_{1}\) 局部 \(\theta\) -可积，此时

\[
g _ {2} \cdot \theta_ {1} = g _ {2} \cdot (g _ {1} \cdot \theta) = (g _ {2} g _ {1}) \cdot \theta\tag{10}
\]

（`结合律公式'）。

依命题 4 的推论，说 \(g_{2}\) 是 \(\theta_{1}\) -可测的等价于说 \(g_{2}g_{1}\) 是 \(\theta\) -可测的。设这个条件成立。对每个函数 \(f \in \mathcal{K}_{+}(T)\) ，我们依命题 2 和命题 3 有

\[
\int^ {\bullet} | g _ {2} | f d | \theta_ {1} | = \int^ {\bullet} | g _ {2} | f | g _ {1} | d | \theta | = \int^ {\bullet} | g _ {2} g _ {1} | f d | \theta |.
\]

因此，说 \(g_{2}\) 局部 \(\theta_{1}\) -可积等价于说 \(g_{2}g_{1}\) 局部 \(\theta\) -可积。设这个条件成立，依定理 1 我们有

\[
\int f d \left(g _ {2} \cdot \theta_ {1}\right) = \int f g _ {2} d \theta_ {1} = \int f g _ {2} g _ {1} d \theta = \int f d \left(g _ {2} g _ {1} \cdot \theta\right),
\]

这个公式等价于 (10)。
% semantic-fragments: legacy-input-seam
\relax{}
